\section{Methods}
\label{sec:methods}

\subsection{Method I: Distributional few-shot personalization (\dfsp)}
\label{sec:dfsp}

\paragraph{Decomposition.}
We write the human target as the judge's prediction plus a discrepancy,
\[
  p_H(z)=p_J(z)+\delta(z).
\]
Here $p_J(z)$ is the first coordinate of $z$. If the judge is informative, $\delta$ is simpler than
$p_H$: it is smaller, or smoother, or both. It can then be estimated from far fewer labels than
$p_H$ itself.

\paragraph{Local smoothing.}
Following \citet{li2026personalizing}, for $\theta=(\theta_1,\theta_2)$ with $\theta_1\ge0$ and
$0<\theta_2\le1$, and a center $z$, define the truncated offset
\[
  \omega_{\theta,z}(z')=p_J(z)+\min\bigl\{|p_J(z')-p_J(z)|,\ \theta_1\norm{z'-z}^{\theta_2}\bigr\}\,
  \operatorname{sgn}\{p_J(z')-p_J(z)\}.
\]
This operation limits how much the offset may vary around $z$, and $\theta$ acts as a borrowing
dial. When $\theta_1=0$, $\omega_{\theta,z}\equiv p_J(z)$, and the estimator below reduces to a purely
human kernel estimate. When $\theta_1$ is large, the full offset is used.

\paragraph{Estimator.}
Let $\cI_n$ index the revealed labels and let $K_h(z',z)=\ind\{\norm{z'-z}_\infty\le h\}$. Define
\[
  \hat\delta_{\theta,h}(z)=\frac{\sum_{i\in\cI_n}K_h(Z_i,z)\{Y_i-\omega_{\theta,z}(Z_i)\}}
                                 {1\vee\sum_{i\in\cI_n}K_h(Z_i,z)},
  \qquad
  \hat p_{\theta,h}(z)=\Pi_{[0,1]}\{p_J(z)+\hat\delta_{\theta,h}(z)\},
\]
where $\Pi_{[0,1]}$ clips to $[0,1]$. The pair $(\theta,h)$ is chosen from a grid by validation
Brier score on a held-out part of the revealed labels. Because the grid contains $\theta_1=0$, a
judge that does not help is automatically switched off.

\paragraph{Scalar versus distributional localization.}
Scalar FSP (\sfsp) localizes on $p_J$ alone. \dfsp{} localizes on $Z=(p_J,t_J,o_J)$. The two
procedures target different quantities, $\E(Y\mid p_J)$ and $\E(Y\mid Z)$. The comparison between
them decomposes exactly as
\[
  R_X(\hat p_{\mathrm S})-R_X(\hat p_{\mathrm D})
  =G-\{R(\hat p_{\mathrm D})-R(\hat p_{\mathrm S})\},
  \qquad
  G=\E\{\E(Y\mid Z)-\E(Y\mid p_J)\}^2 .
\]
Here $R_X$ is the error against the per-comparison probability $P(Y=1\mid X)$. Each $R(\cdot)$ is
measured against its own target, so the braced term is the extra estimation cost of localizing in
three dimensions instead of one. The judge distribution helps exactly when its information gain
$G$ exceeds that extra cost. $G$ is a variable-importance quantity
\citep{williamson2021nonparametric,williamson2023general}, and we test it directly
(Section~\ref{sec:experiments}).

\subsection{Method II: Choosing which comparisons to label}
\label{sec:acquisition}

We choose the $n$ comparisons to label from the judged pool. Let $q$ be the density of $Z$ in the
pool, $\sigma^2(z)=p_H(z)\{1-p_H(z)\}$ the human label noise, and $\gamma_1(z)$ the local
H\"older constant of the discrepancy $\delta$. Suppose labels are drawn with density $p$ and the
kernel uses a locally chosen bandwidth. Minimizing the resulting bias--variance upper bound on the
risk over $p$ gives the design
\[
  p^{*}(z)\;\propto\;q(z)^{\frac{2\gamma_2+d}{4\gamma_2+d}}\;
  \gamma_1(z)^{\frac{2d}{4\gamma_2+d}}\;\sigma(z)^{\frac{4\gamma_2}{4\gamma_2+d}} .
\]
In words, label comparisons that are \emph{common}, on which \emph{humans are split}, and where
the judge-to-human discrepancy \emph{varies}. A constant bias of the judge is cheap to learn
however large it is. Labels should not simply go where the judge is unsure.

In practice we proceed as follows:
\begin{enumerate}
  \item Reveal a small uniform pilot sample.
  \item Fit \dfsp{} on the pilot sample.
  \item Estimate $\sigma$ from the pilot fit.
  \item Estimate $\gamma_1$ from local differences of the fitted discrepancy, at a scale larger than
        the pilot bandwidth.
  \item Plug the estimates into $p^{*}$, mix the result with $q$ for robustness, and sample the
        remaining labels from the pool accordingly.
\end{enumerate}
For H\"older classes no design can improve the minimax \emph{rate} \citep{castro2005faster}. The
goal of this rule is a better constant, that is, fewer labels for the same accuracy.

\subsection{Method III: Post-training against the personalized target (\fsppo)}
\label{sec:post-training}

Let $\pi_{\mathrm{ref}}$ be the initial model and $\pi_\vartheta$ the model being trained. Following
DPO \citep{rafailov2023dpo}, define
\[
  d_\vartheta(X)=\beta\Bigl[\log\frac{\pi_\vartheta(y_A\mid x)}{\pi_{\mathrm{ref}}(y_A\mid x)}
  -\log\frac{\pi_\vartheta(y_B\mid x)}{\pi_{\mathrm{ref}}(y_B\mid x)}\Bigr].
\]
We train on all judged pairs with the soft-label objective
\[
  \widehat R(\vartheta;\hat p)=-\frac1N\sum_{i=1}^{N}
  \Bigl[\hat p_i\log\sigma\{d_\vartheta(X_i)\}+(1-\hat p_i)\log\sigma\{-d_\vartheta(X_i)\}\Bigr],
\]
where $\hat p_i=\hat p(Z_i)$ comes from \dfsp. The loss itself is standard
\citep{lee2024rlaif,furuta2024gapo}. What is new is the target $\hat p$ and the guarantee it carries.

Let $R(\vartheta;p)$ denote the population version of this objective, and let $p_X=P(Y=1\mid X)$.

\begin{proposition}[From estimation error to post-training error]
\label{prop:bridge}
Suppose $|d_\vartheta|\le D$, and $\hat\vartheta$ minimizes $R(\cdot;\hat p)$ up to an error
$\varepsilon$ (optimization plus generalization). Then
\[
  R(\hat\vartheta;p_X)-\inf_\vartheta R(\vartheta;p_X)\;\le\;\varepsilon+2D\,R_X(\hat p)^{1/2}.
\]
Suppose in addition the class $\{d_\vartheta\}$ is convex. Then the bound improves to
$2\varepsilon+2R_X(\hat p)/c_D$, where $c_D=\sigma(D)\{1-\sigma(D)\}$.
\end{proposition}

\noindent\emph{Proof idea.} The objective is linear in the target:
$\ell(d;p)-\ell(d;p')=-(p-p')\,d$. Replacing $p_X$ by $\hat p$ therefore shifts every candidate's
risk by at most $D\,\E|\hat p-p_X|$, which gives the first bound. For a convex class, the logistic
loss grows quadratically around its minimizer, which gives the second.

The proposition turns the human-label cost of estimating $p_H$ into a bound on the policy's excess
preference risk. $R_X$ includes the part of human preference not explained by $Z$, which no
amount of human labels removes.
